% vim: set textwidth=78 autoindent:
% !TeX root = user_guide.tex

\chapter{Externe Python Plugins verwenden}\label{sec:external_plugins}
\index{Plugins!Externe Plugins}

% when the revision of a chapter has been finalized, 
% comment out the following line:
% \updatedisclaimer

Externe QGIS Plugins sind in der Sprache Python geschrieben. Sie sind auf
einem Server entweder im sogenannten 'Official' oder 'User contributed'
Repository abgelegt. Dar�berhinaus gibt es noch zahlreiche, weitere externe
Repositories, die von den jeweiligen Autoren gepflegt werden. 

Tabelle \ref{tab:external_plugins} zeigt
eine Liste der aktuellen Plugins im 'Official' Repository mit einer kurzen
Beschreibung. Eine detailierte Dokumentation zur Anwendung, die notwendige
QGIS Version, Heimatseite des Plugins, Autoren und weitere wichtige
Informationen werden mit dem jeweiligen Plugin bereitgestellt und sind nicht
Teil dieses Handbuches. \footnote{Auch Updates von Kern-Plugins k�nnen im
'Official' Repository bereitgestellt werden.}
\footnote{Python Plugin Installer, fTools und Mapserver Export sind auch
Python Plugins. Sie sind gleichzeitig aber auch Teil des QGIS Quellcodes und
damit automatisch wie ein Kern-Plugin im Plugin Manager ausw�hlbar (siehe
Kapitel~\ref{sec:load_external_plugin}).}

Eine aktuelle Liste der moderierten externen Plugins finden Sie auch im QGIS
Official Repository unter der URL:
\url{http://www.qgis.org/de/herunterladen/plugins-de.html}, sowie �ber den
Python Plugin Installer im Men� \dropmenuopttwo{plugin_installer}{Python
Plugins herunterladen} in QGIS selbst. 

\begin{table}[H]
\centering
\caption{Externe QGIS Python Plugins im Official Repository}
\label{tab:external_plugins}\medskip
\small
 \begin{tabular}{|l|l|p{8cm}|}
\hline \textbf{Icon} & \textbf{Externes Plugin} & \textbf{Beschreibung}\\
\hline
\includegraphics[width=0.7cm]{zoom2point_icon}
 & Zoom To Point \index{Plugins!Zoom To Point} & Zoomt zu einer Koordinate,
die im Dialog angegeben wurde. Ein Zoom-Level kann zur Kontrolle des
Kartenausschnitts angegeben werden. \\
\hline
\includegraphics[width=0.7cm]{plugin_installer}
& Plugin Installer \index{Plugins!Python Plugin Installer} & Der aktuelle
Python Plugin Installer. \\
\hline
\end{tabular}
\end{table}

Weitere Informationen zur Installation von externen Python Plugins finden Sie
in Kapitel~\ref{sec:load_external_plugin}.

\begin{Tip} \caption{\textsc{Weitere externe Plugin Repositorys hinzuf�gen}}
Um ein 'User contributed' oder Anwender-Repository hinzuzuf�gen, �ffnen
Sie den Plugin Installer im Men� \mainmenuopt{Plugins} \arrow
\dropmenuopttwo{plugin_installer}{Python Plugins herunterladen}. Nun wechseln
Sie in den Reiter \tab{Repositorys} und klicken auf den Knopf \button{Plugin
Repositorys hinzuf�gen}. Sie k�nnen bereits hinzugef�gte Plugin Repositorys
auch deaktivieren oder ganz aus der Liste l�schen.
\end{Tip}

\FloatBarrier
